\begin{frame}[t]{\ev{\tagO}More attempts found more solutions}
\vspace{.1cm}
\lead{SWE-bench Lite, issues solved by at least one attempt (DeepSeek-Coder-V2-Instruct).}
\begin{columns}[T,onlytextwidth]
\column{.42\textwidth}
\begin{tabular}{@{}lr@{}}
\toprule
attempts per issue & issues solved\\
\midrule
1 sample & 15.9\%\\
250 samples & \textbf{56\%}\\
\midrule
gain & \textbf{40.1 points}\\
\bottomrule
\end{tabular}
\column{.54\textwidth}
{\small
\begin{ticks}
\item[\yes] This measures coverage: some attempt in the collection is correct.
\item[\no] It does not measure selection of that attempt by a production verifier.
\item[\no] With an imperfect verifier and a cost for false acceptance, more sampling can reduce utility. The optimum is often fewer than ten attempts.
\end{ticks}\par}
\end{columns}
\claim{The remaining problem is to recognise the correct candidate.}
\sources{\href{https://arxiv.org/abs/2407.21787}{Brown et al., Large Language Monkeys, arXiv:2407.21787}, abstract. \href{https://arxiv.org/abs/2411.17501}{Stroebl, Kapoor \& Narayanan, The Limits of Inference Scaling Through Resampling, arXiv:2411.17501}, abstract: when false positives carry negative utility, the optimal number of attempts is often fewer than 10.}
\note{[0:45] Repeated sampling raises coverage. Brown and colleagues sampled DeepSeek-Coder-V2-Instruct on SWE-bench Lite. One sample solved 15.9 percent of issues. With 250 samples, at least one attempt solved 56 percent. This measures the coverage of the candidate collection. A production verifier must still pick the correct candidate. Stroebl and colleagues show that an imperfect verifier limits the gain. When false acceptance has a cost, the best number of attempts is often below ten. Our loop already had checks and reviews.}
\end{frame}
